% backmatter.tex — index, glossary of names, sources

\chapter*{A note on the text}
\addcontentsline{toc}{chapter}{A note on the text}
\markboth{A note on the text}{A note on the text}

\noindent
The Latin texts translated in this book are drawn from the Perseus Digital
Library and, where Perseus does not yet hold a work, from the Latin Library.
The order of the works follows the date of their composition or delivery as
attested by ancient sources and the standard scholarly chronologies. Where
dating is disputed --- which is the case for many of the letters and a few
of the early speeches --- the choice has been recorded in the entry's
metadata in the project repository (\texttt{meta/works.yaml}), and the
contested date is flagged in the headnote.

Section numbers in the margin are those of the standard scholarly editions
(Teubner and Oxford Classical Texts) and are how the works are cited in the
secondary literature.

% Profile-specific apparatus (glossary of names, Greek catalogue, letter
% network, cross-references, allusions, translator's notes) is generated
% from data/ sidecars by scripts/generate_backmatter.py. See
% build/PROFILES.md.
%
% In the reading profile (default) the generated file is empty. In study and
% scholar profiles it carries the additional chapters.
\IfFileExists{backmatter-generated.tex}{\input{backmatter-generated.tex}}{%
  \chapter*{Glossary of names}
  \addcontentsline{toc}{chapter}{Glossary of names}
  \markboth{Glossary of names}{Glossary of names}
  \noindent
  \textit{To follow.} The glossary, Greek catalogue, and other apparatus
  are generated from \texttt{data/} sidecars by
  \texttt{scripts/generate\_backmatter.py}. Run that script (or
  \texttt{scripts/assemble\_book.py --profile study}) to populate this
  section.
}
